\section{Conclusion}
\label{sec:conclusion}

SecretHandle moves secret use for mobile AI agents out of the planner's observation and into an Android reference monitor. The OS mints a typed opaque handle where a secret first arrives, intersects source and task constraints, authenticates the record, and reveals only a non-semantic capability. \texttt{ENTER(handle)} succeeds only for a live task and trusted target that satisfy destination policy; atomic lifecycle state limits replay; and principal-specific redaction prevents the next screenshot or accessibility tree from simply revealing the entered value.

The formal analysis identifies what the design can and cannot promise. Given an intact reference monitor, authentic target identity, complete modeled-channel redaction, value-oblivious public behavior, and a secure authenticator, the construction provides conditional non-disclosure, destination confinement, unforgeability, and single-use safety. It does not prevent misuse within an allowed destination, compromise of the OS, or leakage through an omitted channel.

The next scientific step is implementation, not stronger prose. The AOSP prototype must demonstrate source integration, WebView and custom-view behavior, atomic recovery, observation closure, and acceptable overhead on real login, 2FA, and checkout flows. Until those measurements exist, the empirical contribution remains a reproducible protocol and the proposed trends remain placeholders.
